\section{Read results at the caller boundary}
At \texttt{filter\_ready}, the mixer-style caller has returned from the filter and
output memory is complete. With the default impulse and count sixteen,
\texttt{v0}=16 and \texttt{v1}=0. The words \texttt{processed} and \texttt{clipped}
still contain 99 because their stores follow this checkpoint. At
\texttt{results\_ready}, the words contain sixteen and zero.

The callable checks vary input, coefficient and destination byte offsets and exchange
the two large arenas' input/output roles. They verify all output bytes and every
caller-data access, saved registers, supplied return continuations and fresh result
initialization. Separate tests relocate mapped data above four GiB. These tests do
not establish high text or high stack operation. Reads from submission-owned constants
are accepted, and argument-register values after their last use are not constrained.

The construction lesson will inspect particular registers while building this sequence.
Final free-writing grading must accept another implementation with equivalent arithmetic
and ownership. A requirement to demonstrate vectorization needs instructor review of
that strategy in addition to behavioral checks.
